\section{Literal affine coefficients in projective factorizations}

\begin{lemma}[Transport along an actual coefficient-ring isomorphism]
\label{lem:fc-projective-factorization-ring-isomorphism}
\leanok
\lean{EnvelopeProperty.ProjectiveFactorizationRingIso.postcomp_inverse}
\uses{thm:fc-whole-affine-bundle-polynomial-embedding}
Let $e:R\simeq S$ be a ring isomorphism.  If $f:Y\to\operatorname{Spec}R$
admits a finite polynomial projective factorization, then its composite
with $\operatorname{Spec}(e^{-1})$ admits such a factorization over the
literal coefficient ring $S$.
\end{lemma}
\begin{proof}
Base change the original polynomial closed embedding along
$\operatorname{Spec}e:\operatorname{Spec}S\simeq\operatorname{Spec}R$.
The base-changed projective space is the original polynomial
projective space over $S$, and the base-changed immersion is closed.
The map from $Y$ to the base change given by the identity of $Y$ and
$f$ followed by $\operatorname{Spec}(e^{-1})$ is an isomorphism, since
its composite with the first projection is the identity and that
first projection is an isomorphism.  Precompose the base-changed
closed embedding with this map.  The second-projection equality of
the pullback gives the required original composite as the structure
morphism.
\end{proof}

\begin{theorem}[A literal polynomial factorization of a projective morphism to a spectrum]
\label{thm:fc-affine-projectivity-polynomial-factorization}
\leanok
\lean{EnvelopeProperty.ProjectiveFactorizationRingIso.of_isProjective}
\uses{cor:fc-projective-bundle-to-polynomial-factorization,lem:fc-projective-factorization-ring-isomorphism}
Every morphism $f:Y\to\operatorname{Spec}R$ that is projective in the
finite quasi-coherent projective-bundle sense admits a closed
factorization through a finite polynomial projective space over the
literal original ring $R$.
\end{theorem}
\begin{proof}
The whole affine-bundle comparison constructs a polynomial
factorization of $f$ followed by the canonical affine isomorphism
$\operatorname{Spec}R\simeq\operatorname{Spec}\Gamma(\operatorname{Spec}R,
\mathcal O)$.  Apply the preceding coefficient transport using the
canonical ring isomorphism
$\Gamma(\operatorname{Spec}R,\mathcal O)\simeq R$.  The inverse
spectrum morphism of this canonical isomorphism is exactly the
inverse of that affine isomorphism.  Their composites cancel, so the
resulting structure morphism is the literal original $f$.
\end{proof}
